\newpage
\begingroup\sloppy
\begin{thebibliography}{114}
\addcontentsline{toc}{section}{References}
\bibitem{r1} I. Gati, M. Krausz, and S. H. Osipow, ``A taxonomy of difficulties in career decision making,'' \textit{Journal of Counseling Psychology}, vol. 43, no. 4, pp. 510-526, 1996.

\bibitem{r2} N. Levin, S. Udayar, Y. Lipshits-Braziler, I. Gati, and J. Rossier, ``The structure of the Career Decision-Making Difficulties Questionnaire across 13 countries,'' \textit{Journal of Career Assessment}, vol. 31, no. 1, pp. 129-148, 2023, doi: 10.1177/10690727221099226.

\bibitem{r3} T. M. Trinh, T. T. K. Le, K. M. A. Le, C. Nguyen, and T. N. Tran, ``Shaping choices: Factors influencing Vietnamese high school students' transition to higher education,'' 2024. [Verify venue/pages before citation.]

\bibitem{r4} ``Nearly 20\% of university students lack career direction, survey shows,'' \textit{VietnamNet}, May 23, 2026. [Grey literature; survey of \textasciitilde{}9,200 students across 10 Vietnamese universities, reported at the ``Building Happy Universities in Vietnam'' seminar.]

\bibitem{r5} T. Q. Tran, N. B. T. Vu, D. Van Le, and L. H. Vu, ``Vertical and horizontal job-education mismatches and wages: A quantitative analysis of university graduates in Vietnam,'' \textit{International Journal of Educational Development}, vol. 117, art. 103323, 2025, doi: 10.1016/j.ijedudev.2025.103323.

\bibitem{r6} D. J. Pittenger, ``Cautionary comments regarding the Myers-Briggs Type Indicator,'' \textit{Consulting Psychology Journal: Practice and Research}, vol. 57, no. 3, pp. 210-221, 2005.

\bibitem{r7} The Myers \& Briggs Foundation, ``MBTI reliability and validity,'' myersbriggs.org, 2023. [Grey literature; publisher-affiliated source reporting MBTI Manual data.]

\bibitem{r8} J. M. Phillips, ``Effects of realistic job previews on multiple organizational outcomes: A meta-analysis,'' \textit{Academy of Management Journal}, vol. 41, no. 6, pp. 673-690, 1998, doi: 10.2307/256964.

\bibitem{r9} D. R. Earnest, D. G. Allen, and R. S. Landis, ``Mechanisms linking realistic job previews with turnover: A meta-analytic path analysis,'' \textit{Personnel Psychology}, vol. 64, no. 4, pp. 865-897, 2011, doi: 10.1111/j.1744-6570.2011.01230.x.

\bibitem{r10} R. W. Lent, S. D. Brown, and G. Hackett, ``Toward a unifying social cognitive theory of career and academic interest, choice, and performance,'' \textit{Journal of Vocational Behavior}, vol. 45, no. 1, pp. 79-122, 1994.

\bibitem{r11} A. Byars-Winston, J. Diestelmann, J. N. Savoy, and W. T. Hoyt, ``Unique effects and moderators of effects of sources on self-efficacy: A model-based meta-analysis,'' \textit{Journal of Counseling Psychology}, vol. 64, no. 6, pp. 645-658, 2017, doi: 10.1037/cou0000219.

\bibitem{r12} D. A. Kolb, \textit{Experiential Learning: Experience as the Source of Learning and Development}, 2nd ed. Upper Saddle River, NJ, USA: Pearson Education, 2015 (orig. 1984).

\bibitem{r13} J. P. Wanous, \textit{Organizational Entry: Recruitment, Selection, Orientation, and Socialization of Newcomers}, 2nd ed. Reading, MA, USA: Addison-Wesley, 1992.

\bibitem{r14} S. L. Premack and J. P. Wanous, ``A meta-analysis of realistic job preview experiments,'' \textit{Journal of Applied Psychology}, vol. 70, no. 4, pp. 706-719, 1985.

\bibitem{r15} G. M. McEvoy and W. F. Cascio, ``Strategies for reducing employee turnover: A meta-analysis,'' \textit{Journal of Applied Psychology}, vol. 70, no. 2, pp. 342-353, 1985.

\bibitem{r16} M. Graef, ``Realistic job previews,'' Quality Improvement Center for Workforce Development, University of Nebraska--Lincoln, Feb. 2020. [Grey literature; umbrella review.]

\bibitem{r17} D. A. Cook et al., ``Technology-enhanced simulation for health professions education: A systematic review and meta-analysis,'' \textit{JAMA}, vol. 306, no. 9, pp. 978-988, 2011, doi: 10.1001/jama.2011.1234.

\bibitem{r18} W. C. McGaghie, S. B. Issenberg, E. R. Cohen, J. H. Barsuk, and D. B. Wayne, ``Does simulation-based medical education with deliberate practice yield better results than traditional clinical education? A meta-analytic comparative review of the evidence,'' \textit{Academic Medicine}, vol. 86, no. 6, pp. 706-711, 2011.

\bibitem{r19} O. Chernikova, N. Heitzmann, M. Stadler, D. Holzberger, T. Seidel, and F. Fischer, ``Simulation-based learning in higher education: A meta-analysis,'' \textit{Review of Educational Research}, vol. 90, no. 4, pp. 499-541, 2020, doi: 10.3102/0034654320933544.

\bibitem{r20} V. J. Shute, ``Stealth assessment in computer-based games to support learning,'' in \textit{Computer Games and Instruction}, S. Tobias and J. D. Fletcher, Eds. Charlotte, NC, USA: Information Age Publishing, 2011, pp. 503-524.

\bibitem{r21} V. J. Shute, L. Wang, S. Greiff, W. Zhao, and G. Moore, ``Measuring problem solving skills via stealth assessment in an engaging video game,'' \textit{Computers in Human Behavior}, vol. 63, pp. 106-117, 2016, doi: 10.1016/j.chb.2016.05.047.

\bibitem{r22} F. Delamare Le Deist and J. Winterton, ``What is competence?'' \textit{Human Resource Development International}, vol. 8, no. 1, pp. 27-46, 2005, doi: 10.1080/1367886042000338227.

\bibitem{r23} S. J. Motowidlo, M. D. Dunnette, and G. W. Carter, ``An alternative selection procedure: The low-fidelity simulation,'' \textit{Journal of Applied Psychology}, vol. 75, no. 6, pp. 640-647, 1990.

\bibitem{r24} M. A. McDaniel, F. P. Morgeson, E. B. Finnegan, M. A. Campion, and E. P. Braverman, ``Use of situational judgment tests to predict job performance: A clarification of the literature,'' \textit{Journal of Applied Psychology}, vol. 86, no. 4, pp. 730-740, 2001, doi: 10.1037/0021-9010.86.4.730.

\bibitem{r25} M. S. Christian, B. D. Edwards, and J. C. Bradley, ``Situational judgment tests: Constructs assessed and a meta-analysis of their criterion-related validities,'' \textit{Personnel Psychology}, vol. 63, no. 1, pp. 83-117, 2010.

\bibitem{r26} J. C. Flanagan, ``The critical incident technique,'' \textit{Psychological Bulletin}, vol. 51, no. 4, pp. 327-358, 1954.

\bibitem{r27} D. L. Whetzel, T. S. Sullivan, and R. A. McCloy, ``Situational judgment tests: An overview of development practices and psychometric characteristics,'' \textit{Personnel Assessment and Decisions}, vol. 6, no. 1, 2020.

\bibitem{r28} M. A. McDaniel, N. S. Hartman, D. L. Whetzel, and W. L. Grubb, ``Situational judgment tests, response instructions, and validity: A meta-analysis,'' \textit{Personnel Psychology}, vol. 60, no. 1, pp. 63-91, 2007, doi: 10.1111/j.1744-6570.2007.00065.x.

\bibitem{r29} J. Hamari, J. Koivisto, and H. Sarsa, ``Does gamification work? --- A literature review of empirical studies on gamification,'' in \textit{Proc. 47th Hawaii Int. Conf. System Sciences (HICSS)}, 2014, pp. 3025-3034, doi: 10.1109/HICSS.2014.377.

\bibitem{r30} M. Sailer and L. Homner, ``The gamification of learning: A meta-analysis,'' \textit{Educational Psychology Review}, vol. 32, no. 1, pp. 77-112, 2020, doi: 10.1007/s10648-019-09498-w.

\bibitem{r31} S. Bai, K. F. Hew, and B. Huang, ``Does gamification improve student learning outcome? Evidence from a meta-analysis and synthesis of qualitative data in educational contexts,'' \textit{Educational Research Review}, vol. 30, art. 100322, 2020, doi: 10.1016/j.edurev.2020.100322.

\bibitem{r32} D. B. Clark, E. E. Tanner-Smith, and S. S. Killingsworth, ``Digital games, design, and learning: A systematic review and meta-analysis,'' \textit{Review of Educational Research}, vol. 86, no. 1, pp. 79-122, 2016.

\bibitem{r33} L. Li, K. F. Hew, and J. Du, ``Gamification enhances student intrinsic motivation, perceptions of autonomy and relatedness, but minimal impact on competency: A meta-analysis and systematic review,'' \textit{Educational Technology Research and Development}, vol. 72, no. 2, pp. 765-796, 2024, doi: 10.1007/s11423-023-10337-7.

\bibitem{r34} Y. Zeng et al., ``The effects of gamification on academic performance: A meta-analysis of 2008-2023,'' \textit{British Journal of Educational Technology}, 2024.

\bibitem{r35} ``Effectiveness of the role-play method: A meta-analysis,'' \textit{International Journal of Instruction}, vol. 18, no. 1, 2025.

\bibitem{r36} T. Sitzmann, ``A meta-analytic examination of the instructional effectiveness of computer-based simulation games,'' \textit{Personnel Psychology}, vol. 64, no. 2, pp. 489-528, 2011, doi: 10.1111/j.1744-6570.2011.01190.x.

\bibitem{r37} P. Wouters, C. van Nimwegen, H. van Oostendorp, and E. D. van der Spek, ``A meta-analysis of the cognitive and motivational effects of serious games,'' \textit{Journal of Educational Psychology}, vol. 105, no. 2, pp. 249-265, 2013.

\bibitem{r38} R. Breien and F. Fossøy, ``Narrative categorization in digital game-based learning: Engagement, motivation \& learning,'' \textit{British Journal of Educational Technology}, vol. 52, no. 1, pp. 91-111, 2021, doi: 10.1111/bjet.13004.

\bibitem{r39} ``Randomized reward mechanisms (loot boxes) and their association with problem gambling and wellbeing: A pre-registered study,'' \textit{Addictive Behaviors} / PMC10731324, 2023. [Verify authors/venue before citation.]

\bibitem{r40} R. M. Ryan, C. S. Rigby, and A. Przybylski, ``The motivational pull of video games: A self-determination theory approach,'' \textit{Motivation and Emotion}, vol. 30, no. 4, pp. 344-360, 2006, doi: 10.1007/s11031-006-9051-8.

\bibitem{r41} J. L. Holland, \textit{Making Vocational Choices: A Theory of Vocational Personalities and Work Environments}, 3rd ed. Odessa, FL, USA: Psychological Assessment Resources, 1997.

\bibitem{r42} H.-B. Sheu et al., ``Testing the choice model of social cognitive career theory across Holland themes: A meta-analytic path analysis,'' \textit{Journal of Vocational Behavior}, vol. 76, no. 2, pp. 252-264, 2010.

\bibitem{r43} D. E. Super, ``A life-span, life-space approach to career development,'' \textit{Journal of Vocational Behavior}, vol. 16, no. 3, pp. 282-298, 1980.

\bibitem{r44} J. D. Krumboltz, ``The happenstance learning theory,'' \textit{Journal of Career Assessment}, vol. 17, no. 2, pp. 135-154, 2009.

\bibitem{r45} A. Bandura, \textit{Self-Efficacy: The Exercise of Control}. New York, NY, USA: W. H. Freeman, 1997.

\bibitem{r46} R. W. Lent, G. W. Ireland, L. T. Penn, T. R. Morris, and R. Sappington, ``Sources of self-efficacy and outcome expectations for career exploration and decision-making: A test of the social cognitive model of career self-management,'' \textit{Journal of Vocational Behavior}, vol. 99, pp. 107-117, 2017, doi: 10.1016/j.jvb.2017.01.002.

\bibitem{r47} B. Chang, X. Zhong, and Z. He, ``Sources of self-efficacy in career contexts: A longitudinal perspective,'' \textit{Current Psychology}, vol. 44, pp. 7393-7403, 2025, doi: 10.1007/s12144-025-07699-x.

\bibitem{r48} G. Kurdi, J. Leo, B. Parsia, U. Sattler, and S. Al-Emari, ``A systematic review of automatic question generation for educational purposes,'' \textit{International Journal of Artificial Intelligence in Education}, vol. 30, no. 1, pp. 121-204, 2020, doi: 10.1007/s40593-019-00186-y.

\bibitem{r49} S. Sarsa, P. Denny, A. Hellas, and J. Leinonen, ``Automatic generation of programming exercises and code explanations using large language models,'' in \textit{Proc. 2022 ACM Conf. Int. Computing Education Research (ICER)}, vol. 1, 2022, pp. 27-43, doi: 10.1145/3501385.3543957.

\bibitem{r50} J. Doughty et al., ``Automated educational question generation at different Bloom's skill levels using large language models,'' in \textit{Artificial Intelligence in Education (AIED 2024)}, LNCS, Cham, Switzerland: Springer, 2024, doi: 10.1007/978-3-031-64299-9\_12.

\bibitem{r51} ``Evaluating the instrumental quality of LLM-generated assessment items,'' \textit{Frontiers in Education}, 2026, doi: 10.3389/feduc.2026.1837523.

\bibitem{r52} L. Yan et al., ``Practical and ethical challenges of large language models in education: A systematic scoping review,'' \textit{British Journal of Educational Technology}, vol. 55, no. 1, pp. 90-112, 2024, doi: 10.1111/bjet.13370.

\bibitem{r53} J. S. Park, J. C. O'Brien, C. J. Cai, M. R. Morris, P. Liang, and M. S. Bernstein, ``Generative agents: Interactive simulacra of human behavior,'' in \textit{Proc. 36th Annu. ACM Symp. User Interface Software and Technology (UIST)}, 2023, Art. 2, pp. 1-22, doi: 10.1145/3586183.3606763.

\bibitem{r54} A. Zhu, L. Martin, A. Head, and C. Callison-Burch, ``CALYPSO: LLMs as Dungeon Masters' assistants,'' in \textit{Proc. AAAI Conf. Artificial Intelligence and Interactive Digital Entertainment (AIIDE)}, vol. 19, no. 1, 2023, pp. 380-390, doi: 10.1609/aiide.v19i1.27534.

\bibitem{r55} C. Callison-Burch et al., ``Dungeons and Dragons as a dialog challenge for artificial intelligence,'' in \textit{Proc. 2022 Conf. Empirical Methods in Natural Language Processing (EMNLP)}, 2022, pp. 9379-9393.

\bibitem{r56} ``SNAP: A plan-driven framework for controllable interactive narrative generation,'' arXiv:2601.11529, 2026. [Preprint.]

\bibitem{r57} L. Zheng et al., ``Judging LLM-as-a-judge with MT-Bench and Chatbot Arena,'' in \textit{Advances in Neural Information Processing Systems (NeurIPS) 36, Datasets and Benchmarks Track}, 2023, pp. 46595-46623.

\bibitem{r58} Y. Liu, D. Iter, Y. Xu, S. Wang, R. Xu, and C. Zhu, ``G-Eval: NLG evaluation using GPT-4 with better human alignment,'' in \textit{Proc. 2023 Conf. Empirical Methods in Natural Language Processing (EMNLP)}, 2023, pp. 2511-2522, doi: 10.18653/v1/2023.emnlp-main.153.

\bibitem{r59} ``Reliability without validity: A systematic, large-scale evaluation of LLM-as-a-judge models across agreement, consistency, and bias,'' arXiv:2606.19544, 2026. [Preprint.]

\bibitem{r60} S. Tan, S. Zhuang, K. Montgomery et al., ``JudgeBench: A benchmark for evaluating LLM-based judges,'' in \textit{Proc. Int. Conf. Learning Representations (ICLR)}, 2025. arXiv:2410.12784.

\bibitem{r61} ``CODE-GEN: A human-in-the-loop RAG-based agentic AI system for multiple-choice question generation,'' arXiv:2604.03926, 2026. [Preprint.]

\bibitem{r62} U.S. Department of Labor, ``O*NET Resource Center,'' onetcenter.org. [Government data resource; taxonomy figures per O*NET 28.x, accessed 2026.]

\bibitem{r63} European Commission, ``ESCO --- European Skills, Competences, Qualifications and Occupations,'' esco.ec.europa.eu. [EU data resource; 3,039 occupations / 13,939 skills, ESCO v1.2, accessed 2026.]

\bibitem{r64} N. Lambert, V. Pyatkin, J. Morrison et al., ``RewardBench: Evaluating reward models for language modeling,'' in \textit{Findings of the Association for Computational Linguistics: NAACL 2025}, 2025, pp. 1755-1797. arXiv:2403.13787.

\bibitem{r65} ``EQGBench: Evaluating LLMs' educational question generation,'' arXiv:2508.10005, 2025. [Preprint.]

\bibitem{r66} B. Xu, Y. Bai, Y. Sun et al., ``EduBench: A comprehensive benchmarking dataset for evaluating LLMs in diverse educational scenarios,'' arXiv:2505.16160, 2025. [Preprint.]

\bibitem{r67} M.-A. Côté et al., ``TextWorld: A learning environment for text-based games,'' in \textit{Computer Games (CGW 2018)}, Springer, 2019, pp. 41-75, doi: 10.1007/978-3-030-24337-1\_3.

\bibitem{r68} M. Hausknecht, P. Ammanabrolu, M.-A. Côté, and X. Yuan, ``Interactive fiction games: A colossal adventure,'' in \textit{Proc. AAAI Conf. Artificial Intelligence}, vol. 34, no. 5, 2020, pp. 7903-7910.

\bibitem{r69} M. R. Lynn, ``Determination and quantification of content validity,'' \textit{Nursing Research}, vol. 35, no. 6, pp. 382-385, 1986.

\bibitem{r70} D. F. Polit, C. T. Beck, and S. V. Owen, ``Is the CVI an acceptable indicator of content validity? Appraisal and recommendations,'' \textit{Research in Nursing \& Health}, vol. 30, no. 4, pp. 459-467, 2007, doi: 10.1002/nur.20199.

\bibitem{r71} D. M. Williamson, X. Xi, and F. J. Breyer, ``A framework for evaluation and use of automated scoring,'' \textit{Educational Measurement: Issues and Practice}, vol. 31, no. 1, pp. 2-13, 2012.

\bibitem{r72} K. Krippendorff, \textit{Content Analysis: An Introduction to Its Methodology}, 2nd ed. Thousand Oaks, CA, USA: SAGE, 2004.

\bibitem{r73} A. B. Sai, T. Dixit, D. Y. Sheth, S. Mohan, and M. M. Khapra, ``Perturbation CheckLists for evaluating NLG evaluation metrics,'' in \textit{Proc. 2021 Conf. Empirical Methods in Natural Language Processing (EMNLP)}, 2021, pp. 7219-7234.

\bibitem{r74} P. Verga et al., ``Replacing judges with juries: Evaluating LLM generations with a panel of diverse models,'' arXiv:2404.18796, 2024. [Preprint.]

\bibitem{r75} ``Agentic knowledge tracing: A multi-agent LLM architecture for stealth assessment of financial literacy in serious games,'' arXiv:2606.25358, 2026. [Preprint.]

\bibitem{r76} W. Du, Z. Zhu, X. Xu, H. Che, and S. Chen, ``CareerSim: Gamification design leveraging LLMs for career development reflection,'' in \textit{Extended Abstracts of the 2024 CHI Conf. Human Factors in Computing Systems (CHI EA '24)}, Art. 71, pp. 1-7, doi: 10.1145/3613905.3650928. [Extended abstract.]

\bibitem{r77} P. Pataranutaporn et al., ``Future You: A conversation with an AI-generated future self reduces anxiety, negative emotions, and increases future self-continuity,'' in \textit{Proc. IEEE Frontiers in Education (FIE)}, 2024. arXiv:2405.12514.

\bibitem{r78} H. Jeon et al., ``Letters from Future Self: Augmenting the letter-exchange exercise with LLM-based agents to enhance young adults' career exploration,'' in \textit{Proc. 2025 CHI Conf. Human Factors in Computing Systems (CHI '25)}, Art. 100, pp. 1-21, doi: 10.1145/3706598.3714206.

\bibitem{r79} Z. Wang, Z. Zeng, Y. Li, and Z. Ding, ``CareerPooler: AI-powered metaphorical pool simulation improves experience and outcomes in career exploration,'' arXiv:2509.11461, 2025. [Preprint.]

\bibitem{r80} S. Lee and H. Jeong, ``'I Can Be Anything!' Bridging today and the future through generative AI-driven self-represented career imagination for children,'' in \textit{Proc. 2026 CHI Conf. Human Factors in Computing Systems (CHI '26)}, doi: 10.1145/3772318.3790985.

\bibitem{r81} K. Ozaki, X. Geng, and M. Yamada, ``Promoting career self-regulation in high school students using a career simulation game,'' in \textit{Learning and Collaboration Technologies (HCII 2026)}, LNCS vol. 16733, Cham, Switzerland: Springer, 2026, pp. 416-431, doi: 10.1007/978-3-032-30784-2\_26.

\bibitem{r82} H. Han, B. Park, and K. Seo, ``A self-determination theory-based career counseling chatbot,'' in \textit{Extended Abstracts of the 2025 CHI Conf. Human Factors in Computing Systems (CHI EA '25)}, pp. 1-9, doi: 10.1145/3706599.3720286. [Extended abstract.]

\bibitem{r83} Forage, ``Free virtual job simulations and career prep,'' theforage.com. [Grey literature; accessed 2026, verify before citation.]

\bibitem{r84} EntryLevel, ``AI job simulations,'' entrylevel.net. [Grey literature.]

\bibitem{r85} Anthropos, ``AI job simulations: Verify real skills in work scenarios,'' anthropos.work. [Grey literature.]

\bibitem{r86} careersim.ai, ``AI interview practice and workplace conversation simulator,'' careersim.ai. [Grey literature.]

\bibitem{r87} Latitude, ``AI Dungeon,'' aidungeon.com; see also ``AI Dungeon,'' Wikipedia. [Grey literature.]

\bibitem{r88} Character Technologies, ``Character.AI,'' character.ai. [Grey literature.]

\bibitem{r89} Hidden Door, ``Social roleplaying in your favorite fictional worlds,'' hiddendoor.co; and C. Fink, ``Hidden Door turns fan worlds into licensed, revenue-sharing story platforms,'' \textit{Forbes}, Aug. 14, 2025. [Grey literature.]

\bibitem{r90} ``Ứng dụng tư vấn hướng nghiệp 4.0 JobWay,'' Vietnamese science-and-technology press, vista.gov.vn. [Grey literature.]

\bibitem{r91} JobTest.vn, ``Nền tảng hướng nghiệp và đánh giá năng lực,'' jobtest.vn. [Grey literature.]

\bibitem{r92} Hướng nghiệp Sông An, huongnghiepsongan.com. [Grey literature.]

\bibitem{r93} ``Hướng Nghiệp LwL (Lead with LOF),'' Google Play / fptshop.com.vn. [Grey literature.]

\bibitem{r94} ``Chatbot dự đoán tỉ lệ đỗ đại học và trắc nghiệm hướng nghiệp bằng AI của startup Việt,'' \textit{Tuổi Trẻ}, 2026. [Grey literature.]

\bibitem{r95} Prime Minister of Vietnam, Decision No. 522/QĐ-TTg of 14 May 2018 approving the scheme ``Career guidance and orientation for student streaming in general education, 2018--2025,'' 2018.

\bibitem{r96} ``Tiếp tục đẩy mạnh hướng nghiệp, phân luồng trong giáo dục theo hướng lấy học sinh làm trung tâm,'' \textit{Tạp chí Giáo dục}, 2025. [Grey literature.]

\bibitem{r97} Ho Chi Minh City Department of Education and Training, Official Letter No. 5829/SGDĐT-GDTXNNĐH on career guidance and streaming after upper secondary school, school year 2025--2026, Dec. 2025. [Government document.]

\bibitem{r98} S. J. Hamstra, R. Brydges, R. Hatala, B. Zendejas, and D. A. Cook, ``Reconsidering fidelity in simulation-based training,'' \textit{Academic Medicine}, vol. 89, no. 3, pp. 387-392, 2014, doi: 10.1097/ACM.0000000000000130.

\bibitem{r99} National Assembly of Vietnam, Law No. 91/2025/QH15 on Personal Data Protection, 26 Jun. 2025, in force from 1 Jan. 2026. [Legislation.]

\bibitem{r100} Government of Vietnam, Decree No. 356/2025/NĐ-CP detailing a number of articles of, and measures for implementing, the Law on Personal Data Protection, 31 Dec. 2025, in force from 1 Jan. 2026. [Legislation.]

\bibitem{r101} H. Li et al., ``Agreement between large language models and human raters in essay scoring: A research synthesis,'' arXiv:2512.14561, 2025. [Preprint.]

\bibitem{r102} O. Henkel, L. Hills, A. Boxer, B. Roberts, and Z. Levonian, ``Can large language models make the grade? An empirical study evaluating LLMs' ability to mark short answer questions in K-12 education,'' in \textit{Proc. 11th ACM Conf. Learning @ Scale (L@S)}, 2024, doi: 10.1145/3657604.3664693.

\bibitem{r103} K. P. Yancey, G. Laflair, A. Verardi, and J. Burstein, ``Rating short L2 essays on the CEFR scale with GPT-4,'' in \textit{Proc. 18th Workshop on Innovative Use of NLP for Building Educational Applications (BEA)}, 2023, pp. 576-584, doi: 10.18653/v1/2023.bea-1.49.

\bibitem{r104} M. Stahl, L. Biermann, A. Nehring, and H. Wachsmuth, ``Exploring LLM prompting strategies for joint essay scoring and feedback generation,'' in \textit{Proc. 19th Workshop on Innovative Use of NLP for Building Educational Applications (BEA)}, 2024. arXiv:2404.15845.

\bibitem{r105} V. Hackl, A. E. Müller, M. Granitzer, and M. Sailer, ``Is GPT-4 a reliable rater? Evaluating consistency in GPT-4's text ratings,'' \textit{Frontiers in Education}, vol. 8, 2023, doi: 10.3389/feduc.2023.1272229.

\bibitem{r106} N. Saka, I. Gati, and K. R. Kelly, ``Emotional and personality-related aspects of career-decision-making difficulties,'' \textit{Journal of Career Assessment}, vol. 16, no. 4, pp. 403-424, 2008, doi: 10.1177/1069072708318900.

\bibitem{r107} J. Hacker, A. Carr, M. Abrams, and S. D. Brown, ``Development of the Career Indecision Profile: Factor structure, reliability, and validity,'' \textit{Journal of Career Assessment}, vol. 21, no. 1, pp. 32-41, 2013, doi: 10.1177/1069072712453832.

\bibitem{r108} K. M. Taylor and N. E. Betz, ``Applications of self-efficacy theory to the understanding and treatment of career indecision,'' \textit{Journal of Vocational Behavior}, vol. 22, no. 1, pp. 63-81, 1983, doi: 10.1016/0001-8791(83)90006-4.

\bibitem{r109} J. P. Sampson Jr., G. W. Peterson, J. G. Lenz, R. C. Reardon, and D. E. Saunders, ``The design and use of a measure of dysfunctional career thoughts among adults, college students, and high school students: The Career Thoughts Inventory,'' \textit{Journal of Career Assessment}, vol. 6, no. 2, pp. 115-134, 1998, doi: 10.1177/106907279800600201.

\bibitem{r110} S. H. Osipow, C. G. Carney, J. L. Winer, B. Yanico, and M. Koschier, \textit{The Career Decision Scale}, 3rd ed. Odessa, FL: Psychological Assessment Resources, 1987.

\bibitem{r111} W.-C. Mau, ``Assessing career decision-making difficulties: A cross-cultural study,'' \textit{Journal of Career Assessment}, vol. 9, no. 4, pp. 353-364, 2001, doi: 10.1177/106907270100900403.

\bibitem{r112} P. A. Creed and W. O. Yin, ``Reliability and validity of a Chinese version of the Career Decision-Making Difficulties Questionnaire,'' \textit{International Journal for Educational and Vocational Guidance}, vol. 6, no. 1, pp. 47-63, 2006, doi: 10.1007/s10775-006-0003-3.

\bibitem{r113} L. Sovet, J. Tak, and S. Jung, ``Validation of the Career Decision-Making Difficulties Questionnaire among Korean college students,'' \textit{Journal of Career Assessment}, vol. 23, no. 4, pp. 661-676, 2015, doi: 10.1177/1069072714553556.

\bibitem{r114} J. Cohen, ``A coefficient of agreement for nominal scales,'' \textit{Educational and Psychological Measurement}, vol. 20, no. 1, pp. 37-46, 1960, doi: 10.1177/001316446002000104.
\end{thebibliography}
\endgroup
